Nach dem Gesetz wäre die Kraft von Körper 1 auf Körper 2 und, mit vertauschten Rollen, die Kraft von Körper 2 auf Körper 1
\begin{align}
 F_{12} &= k\cdot A_1\cdot m_2, \qquad F_{21} = k\cdot A_2\cdot m_1
\end{align}
Das Wechselwirkungsprinzip (3.~Newton'sches Axiom) verlangt $F_{12} = F_{21}$ für \emph{alle} Körper, also
\begin{align}
 A_1\,m_2 &= A_2\,m_1 \quad\Leftrightarrow\quad \frac{A_1}{m_1} = \frac{A_2}{m_2}
\end{align}
Das Verhältnis von Oberfläche zu Masse ist aber für verschiedene Körper verschieden: Eine kleine Stahlkugel und ein grosser, leichter Ballon haben ganz unterschiedliche Werte. Die beiden Kräfte wären im Allgemeinen nicht gleich gross.

Das Gesetz ist deshalb $\result{\text{nicht im Einklang}}$ mit den Newton'schen Axiomen. (Beim Newton'schen Gravitationsgesetz $F \propto m_1\,m_2$ ist das Produkt dagegen symmetrisch.)
